\section{Dataset Description}

The reference dataset used in the experiments is the RT-CEA bi-mode dataset
stored in HDF5 format as \texttt{RTCEA\_bimode.hdf5}. It contains
1978 two-dimensional scalar fields representing the density extracted
from Direct Numerical Simulations of the Rayleigh--Taylor instability. Each
raw field is stored as a single-channel array of size \(512 \times 256\) in
single precision. The associated label array has size \(1978 \times 8\);
the eight recorded quantities are the Atwood number, the physical time, the
mixing-zone size, the non-dimensional time, the non-dimensional mixing-zone
size, the relative mixing-zone size, the phase, and the number of initial
modes. In the present chapter, the dataset considered is the bi-mode subset,
so the last label component identifies simulations generated from a two-mode
initial perturbation.

Before being passed to the learning and evaluation pipelines, the fields are
processed in the same way as in the augmentation notebook. First, the raw
arrays are loaded and vertically reversed to recover the plotting convention
used throughout the project. The density field is then normalized and cropped
around the mixing region. Samples for which the mixing zone is not fully
contained inside the cropped domain are discarded. The remaining fields are
resized to \(256 \times 256\), yielding square images compatible with the
image-space and wavelet-space diffusion models used in the experiments.

Because the number of DNS snapshots is modest compared with standard image
generation datasets, a symmetry-based augmentation is applied. For each field,
the code keeps the original sample, applies integer horizontal translations,
and adds the corresponding horizontally mirrored fields. The translations are
implemented with periodic rolls along the horizontal direction, which is
consistent with the lateral periodicity usually assumed for RTI simulations.
The augmentation metadata are appended to the physical labels as two additional
components: the applied horizontal translation and a binary flip indicator.
With the implementation used here, one input field generates \(20\) labelled
fields in total. The augmentation is therefore not intended to create new
physical parameters; it only exploits geometric invariances of the governing
problem to increase the number of training examples while preserving the
underlying RTI statistics.

The notebook \texttt{Augmentation.ipynb} illustrates this pipeline on a small
slice of the dataset, namely snapshots with indices \(100\) to \(119\). This
subset contains \(20\) original fields and becomes \(400\) fields after
augmentation, before the subsequent normalization, cropping, and resizing
steps. This restricted example is used only for visualization and verification
of the preprocessing operations; the full HDF5 file is not loaded repeatedly
when only structural information such as shapes and labels is required.
